% **************************************************
% Macro specifiche per il documento corrente
% **************************************************
% Nome
\newcommand{\docName}{ManualeUtente}
% Nome file
\newcommand{\docFileName}{ManualeUtente.1.0.pdf}
% Versione
\newcommand{\docVers}{1.0}
% Data creazione
\newcommand{\creationDate}{2013-03-07}
% Data ultima modifica
\newcommand{\modificationDate}{2013-03-07}
% Stato in {Approvato, Non approvato}
\newcommand{\docState}{Non approvato}
% Uso in {Interno, Esterno}
\newcommand{\docUsage}{Esterno}
% Destinatari da specificare come nome1\\ &nome2\\ ecc.
\newcommand{\docDistributionList}{Prof. Tullio Vardanega\\ & Prof. Riccardo Cardin}
% Redattori da specificare come nome1\\ &nome2\\ ecc.
\newcommand{\docAuthors}{Elena Zecchinato}
% Approvato da
\newcommand{\approvedBy}{Andrea Rizzi}
% Verificatori
\newcommand{\verifiedBy}{---}
% Perscorso (relativo o assoluto) che punta alla directory contenente shared/
% come sua sottodirectory (per comodità chiamiamola 'doc root').
\newcommand{\docRoot}{..}
% definire se si vuole l'indice delle tabelle
\def\INDICETABELLE{true}
% definire se si vuole l'indice delle figure
\def\INDICEFIGURE{true}

% importa il preambolo condiviso da tutti i documenti
% shared/preamble.tex
%
% Questo documento contiene la parte del preambolo condivisa e viene pertanto
% richiamato nel 'master' di tutti i documenti di progetto.  Al suo interno
% contiene le inclusioni (e le configurazioni) di tutti i package richiesti per
% la compilazione dei documenti, le macro di carattere generale e la definizione
% degli stili di pagina.

\documentclass[a4paper,10pt,openright]{article}

% **************************************************
% Macro generiche
% **************************************************
\newcommand{\team}{Software Synthesis}                    % chi siamo
\newcommand{\email}{software.synthesis@gmail.com}         % e-mail
\newcommand{\caName}{}                                    % titolo capitolato
\newcommand{\caDescr}{}                                   % descrizione
\newcommand{\inglese}[1]{\textit{#1}}

% **************************************************
% Codifica e lingua dei documenti
% **************************************************
\usepackage[utf8x]{inputenc}                              % codifica caratteri dei documenti sorgenti
\usepackage[english,italian]{babel}                       % localizzazione ai fini di sillabazione e cross-references
\usepackage[T1]{fontenc}                                  % codifica font di output

% **************************************************
% Definizione geometria della pagina
% **************************************************
\usepackage[a4paper,head=4cm,top=4.5cm,bottom=3cm,left=3cm,right=3cm,bindingoffset=5mm]{geometry}

% *************************************************
% Intestazioni e piè di pagina personalizzati
% *************************************************
\usepackage{fancyhdr}
% stile normale
\fancypagestyle{normal}{
\fancyhead{}                                              % intestazione
\fancyhead[RE,RO]{
\begin{picture}(0,0)
  \put(-410,0){\includegraphics[width=1.02\textwidth]{header_logo}}
  \put(-410,10){\sffamily\large\leftmark}
\end{picture}
\vspace{-4pt}
}
\renewcommand{\headrulewidth}{.4pt}                       % riga sotto l'intestazione
\cfoot{}                                                  % piè di pagina
\fancyfoot[RO,LE]{\sffamily
  pag.~\thepage{} di \pageref{LastPage}}                  % a dx nelle pag. dispari e a sx in quelle pari
\fancyfoot[RE,LO]{\sffamily\docFileName{} -- v.\docVers}
\renewcommand{\footrulewidth}{.4pt}                       % riga sopra il piè di pagina
}
% stile per gli indici
\fancypagestyle{toc}{
\fancyhead{}                                              % intestazione
\fancyhead[RE,RO]{
\begin{picture}(0,0)
  \put(-410,0){\includegraphics[width=1.02\textwidth]{header_logo}}
\end{picture}
}
\renewcommand{\headrule}{}                                % nessuna riga sotto l'intestazione
\cfoot{}                                                  % piè di pagina
\fancyfoot[RO,LE]{\sffamily\thepage{}}                    % a dx nelle pag. dispari e a sx in quelle pari
\fancyfoot[RE,LO]{\sffamily\docFileName{} -- v.\docVers}
\renewcommand{\footrulewidth}{.4pt}                       % riga sopra il piè di pagina
}

\pagestyle{fancy}                                         % premetto: non so usare bene le marche:
\renewcommand{\sectionmark}[1]{\markboth{#1}{#1}}         % se qualcuno ha idee migliori si faccia avanti!

% **************************************************
% Tabelle
% **************************************************
\usepackage{tabularx}                                     % tabelle di larghezza fissa con una o più colonne variabili
\usepackage{multirow}                                     % colonne con colonne che si estendono per più righe
\usepackage{booktabs}                                     % per inserire l'ambiente table e le righe orizz. nelle tabelle
\usepackage{longtable}			                          % tabelle oltre i limiti di pagina

% **************************************************
% Cross-references e collegamenti ipertestuali
% **************************************************
\usepackage[hidelinks]{hyperref}
\hypersetup{%
  colorlinks=false, linktocpage=false, pdfborder={0,0,0}, pdfstartpage=3, pdfstartview=FitV,%
  urlcolor=Cyan, linkcolor=Cyan, citecolor=Black, %pagecolor=Black,%
  pdftitle={\docName}, pdfauthor={\team}, pdfsubject={}, pdfkeywords={},%
  pdfcreator={pdflatex}, pdfproducer={pdflatex with hyperref package}%
}

% **************************************************
% Immagini e grafica
% **************************************************
\usepackage{graphicx}                                     % supporto ad aspetti avanzati delle immagini
\graphicspath{{\docRoot/pics/}}                           % percorso contenente tutti i file immagini
\usepackage{color}                                        % permette di colorare facilmente il testo

% **************************************************
% Altri pacchetti opzionali
% **************************************************     
\usepackage{lastpage}                                     % per sapere il numero totale di pagine
\usepackage{lipsum}                                       % genera "dummy text" per prove di impaginazione
\usepackage{eurosym}                                      % per il simbolo dell'euro usare \EUR{x} dove x è l'importo


\usepackage[italian]{varioref}

% Fine del preambolo e inizio del documento
\begin{document}

% Inclusione della prima pagina
% shared/firstpage.tex
%
% Questo documento definisce il contenuto della prima pagina, che si suppone
% essere uguale in tutti i documenti.  Oltre al logo e al titolo, la prima
% pagina contiene i metadati relativi al documento in cui viene inclusa.


% rimuove intestazioni e piè di pagina
\pagestyle{empty}

\begin{center}

% logo del gruppo
\includegraphics[width=1.5\textwidth]{logo}

\vspace{1in}

% titolo del documento
{\Huge\bfseries \docName}

\vspace{1in}

% tabella riepilogativa
\begin{tabularx}{.7\textwidth}{>{\bfseries\sffamily}l>{\sffamily}l}
\toprule
\multicolumn{2}{>{\sffamily}c}{Informazioni sul documento}\\
\midrule
Nome file:            & \docFileName\\
Versione:             & \docVers\\
Data creazione:       & \creationDate\\
Data ultima modifica: & \modificationDate\\
Stato:                & \docState\\
Uso:                  & \docUsage\\
Redattori:            & \docAuthors\\
Approvato da:         & \approvedBy\\
Verificatori:         & \verifiedBy\\
\bottomrule
\end{tabularx}

\end{center}

\newpage


% Storico delle modifiche
\section*{Storia delle modifiche}
\begin{longtable}{lp{.3\textwidth}lll}
\toprule
Versione & Descrizione intervento & Redattore & Ruolo & Data\\
\midrule % inserire qui il contenuto della tabella

0.1 & Stesura scheletro documento, ``bozza manuale'' & Elena Zecchinato & Programmatore & 2013-03-07\\
\bottomrule
\end{longtable}
\newpage

% inclusione dell'indice
% shared/toc.tex
%
% Questo file contiene le istruzioni che generano l'indice o gli indici del
% documento (utile nel caso in cui decidessimo di avere anche un indice delle
% tabelle e/o un indice delle figure).

\pagestyle{toc}
\pagenumbering{roman}

\tableofcontents

\newpage



% Alcuni aggiustamenti per le pagine
\pagenumbering{arabic}
\setcounter{page}{1}
\pagestyle{normal}

% Qui ha inizio il documento vero e proprio...
\section{Introduzione}
\subsection{Scopo del prodotto}
\purpose

\subsection{Scopo del documento}
Lo scopo del presente documento è quello di fornire una guida per comprendere e utilizzare al meglio l'applicazione ``\caName''.
Tale documento contiene la descrizione delle principali funzionalità del prodotto ed il modo
per utilizzarle.


\subsection{Destinatario del manuale}
L'applicazione  ``\caName'' è indirizzata agli utenti che desiderano effettuare una comunicazione audio/video tramite un computer in modo semplice  e senza dover installare nessun software esterno.  

\subsection{Requisiti del sistema}
\subsubsection{Requisiti software}
Il corretto funzionamento del software  ``\caName'' è assicurato solo con i seguenti browser:
\begin{itemize}
  \item Google Chrome (versione minima 23.0.1271.97m)
  \item Firefox (versione 17.0 minima)
\end{itemize}

Tuttavia, pur utilizzando tali browser non è garantita la possibilità di controllo delle statistiche. Infatti, tale funzionalità attualmente è disponibile solo per la versione sviluppatori dei browser. Sarà però probabilmente possibile disporne nelle prossime versioni dei browser.
L'utente deve inoltre disporre dei seguenti dispositivi:
\begin{itemize}
  \item microfono e/o  \inglese{webcam}
 \end{itemize}
 
 
\subsubsection{Prerequisiti}
Per poter usufruire di tutte le funzionalità del prodotto è necessario essere autenticati e quindi registrati al sistema.
All'utente non sono richieste conoscenze particolari oltre alla comune esperienza
di navigazione web.

\subsection{Come comunicare il rilevamento di problemi e malfunzionamenti}
% TODO: 
% da decidere come gestire
Nell'eventualità in cui durante l'utilizzo del software \caName{} si dovessero riscontrare malfunzionamenti o comportamenti inattesi che non corrispondono a quanto descritto nel presente manuale, gli utenti sono invitati a inviare le proprie segnalazioni all'indirizzo:
\begin{center}
  \email{}
\end{center}
specificando le seguenti informazioni:
\begin{itemize}[noitemsep,nolistsep]
  \item[-] il browser utilizzato, la versione dello stesso e il sistema operativo;
  \item[-] una descrizione dettagliata del problema riscontrato, le circostanze in cui si è verificato;
  \item[-] i messaggi d'errore che sono stati presentati.
\end{itemize}

\subsection{Come leggere il manuale}
% TODO: 
% da aggiungere alla fine come si legge il manuale

\subsection{Glossario}
\glossaryIntro


\section{Istruzioni per l'uso}

\subsection{Installazione}
L'applicazione non necessita di installazione, è sufficiente aprire il browser e...
% TODO: 
% da aggiungere a che pagine si deve collegare utente
\subsection{Registrazione}
Se l'utente desidera usufruire del servizi offerti dall'applicazione ``\caName'' allora deve effettuare la registrazione al sistema.\\
L'utente può accedere alla pagina registrazione premendo il pulsante Registrati visualizzato in figura......
Per effettuare la registrazione l'utente deve inserire obbligatoriamente uno \inglese{username} e una \inglese{password}.
L'utente per motivi di sicurezza deve inoltre selezionare una domanda segreta tra quelle proposte e fornirne la risposta.
Il nome, il cognome e l'immagine di profilo sono dati facoltativi.

\subsection{Autenticazione}
Ogni volta che l'utente desidera utilizzare i servizi messi a disposizione dall'applicazione deve necessariamente effettuare il \inglese{login} inserendo il suo nome utente e la sua \inglese{password}.
Fiugura....

\subsection{Home Screen dell'applicativo }
pagine principale per spiegazione dei pannelli


\subsection{Chiamata}
Quando l'utente vuole effettuare una chiamata deve selezionare dalla rubrica l'utente che desidera chiamate. A questo punto nel pannello centrale della pagina appaiono i dati dell'utente selezionato, e tramite un \inglese{click} sul pulsante chiama è possibile iniziare la chiamata.
\subsection{Video chiamata}
\subsection{Chat}
\subsection{Impostazioni Account}
\subsection{Segreteria}
\subsection{Storico delle chiamate}
\subsection{Impostazione lingua}
L'utente ha la possibilità di impostare la lingua desiderata premendo il bottone impostaLingua che apre un menu a tendina con la scelta delle possibili lingue.



\end{document}
